\documentclass[first=Haakon,last=Bessert,company=SAP\:SE,location=Dresden,simple,headertitle=KPI-Hub-Integration in fragmentierten Enterprise-Systemlandschaften]{baarticle}

\newacronym{kpi}{KPI}{Key Performance Indicator}
\newacronym{api}{API}{Application Programming Interface}
\newacronym{et}{ETL}{Extract, Transform, Load}
\newacronym{bi}{BI}{Business Intelligence}
\newacronym{mvp}{MVP}{Minimum Viable Product}
\newacronym{rbac}{RBAC}{Role-Based Access Control}
\newacronym{dsr}{DSR}{Design Science Research}
\newacronym{poc}{PoC}{Proof of Concept}

\newglossaryentry{kpihub}{
    name={KPI-Hub},
    description={Zentraler Zugangspunkt („Single Pane of Glass"), der Kennzahlen aus mehreren Quellsystemen integriert oder föderiert bereitstellt, ohne zwingend eine zentrale Datenhaltung vorauszusetzen}
}
\newglossaryentry{datagovernance}{
    name={Data Governance},
    description={Rahmenwerk aus Verantwortlichkeiten, Regeln und Kontrollen für den Umgang mit Daten innerhalb einer Organisation. Beeinflusst welche Integrationsformen zulässig und tragfähig sind}
}
\newglossaryentry{federation}{
    name={Dashboard-Föderation},
    description={Integrationsansatz, bei dem bestehende Dashboards anderer Systeme per Embedding in eine übergeordnete Oberfläche eingebettet werden, ohne die Daten selbst zu kopieren}
}
\newglossaryentry{artefakt}{
    name={Artefakt},
    description={Im Design-Science-Kontext ein durch den Forschungsprozess konstruiertes Objekt (hier: Bewertungsmodell und MVP), das das Problemverständnis illustriert und Lösungspotenziale aufzeigt}
}

\addbibresource{references.bib}

\newcommand{\basimpleabstract}{
    \addcontentsline{toc}{section}{Abstract}
    \thispagestyle{empty}
    Im Unternehmensbereich SAP Signavio Product \& Engineering liegen zentrale Kennzahlen und operative Metriken verteilt über mehrere Systeme und Dashboards vor. Dazu gehören Software- und Plattformlösungen wie z.B. SAP Analytics Cloud, JIRA, ServiceNow sowie PDF- oder Excel-Reports. Diese Fragmentierung erschwert ein konsistentes Gesamtbild um einen Geschäftsbereich effektiv und effizient datenbasiert steuern zu können.

    Für datengestützte Geschäftsbereichsentscheidungen werden aktuell manuell die Kennzahlen und Informationen aus den verschiedenen Software- und Plattformlösungen zusammengetragen und mit Kontext angereicht. Das ist zeitaufwendig, fehleranfällig und erschwert insbesondere die systemübergreifende Einordnung von Ursachen und Auswirkungen, beispielsweise zwischen Ticketlage, Sprint-Performance und operativen \gls{kpi}.

    Ziel dieser Bachelorarbeit ist die Entwicklung eines belastbaren Entscheidungswegs zur Auswahl eines geeigneten \gls{kpihub}-Integrationsansatzes. Unter einem \gls{kpihub} wird dabei ein zentraler Zugangspunkt („single pane of glass”) verstanden, der Kennzahlen aus bestehenden Quellsystemen integriert oder föderiert bereitstellt, ohne eine zentrale Datenhaltung vorauszusetzen. Im Mittelpunkt steht nicht die Frage, welche Lösung technisch grundsätzlich möglich ist, sondern welche Option im Enterprise-Kontext am sinnvollsten ist, wenn technische Machbarkeit, Kosten, Nutzererlebnis, langfristige Wartbarkeit sowie insbesondere Compliance- und Governance-Anforderungen gemeinsam betrachtet werden. Ein zentrales Kriterium ist die rollenbasierte Zugriffskontrolle (\gls{rbac}), damit Nutzerinnen und Nutzer ausschließlich die Informationen und Daten sehen, für die sie berechtigt sind. \gls{datagovernance} wird ergänzend als Rahmen verstanden, der Verantwortlichkeiten, Regeln und Kontrollen für den Umgang mit Daten definiert und damit maßgeblich beeinflusst, welche Integrationsformen zulässig und tragfähig sind.

    Methodisch ist die Arbeit gestaltungsorientiert angelegt und folgt dem \gls{dsr}-Ansatz, bei dem Erkenntnisgewinn wesentlich durch das Konstruieren und Evaluieren eines \gls{artefakt}s entsteht. Ausgehend vom Problem wird ein \gls{artefakt} in Form eines Kriterien- und Bewertungsmodells entwickelt, angewendet und überprüft. Zunächst werden in einer strukturierten Anforderungsanalyse zentrale Stakeholder-Use-Cases erhoben und in priorisierte Anforderungen überführt, einschließlich Erwartungen an Datenaktualität, Performance und Zielgruppen. Parallel erfolgt eine technische Untersuchung der relevanten Quellsysteme mit Blick auf Integrationsmöglichkeiten (u.\,a. \gls{api}s, \gls{et}, Embedding, Authentifizierungs- und Autorisierungsmechanismen) sowie mögliche Einschränkungen durch Lizenzen oder Security-Vorgaben.

    Auf Basis dieser Ergebnisse werden mehrere Lösungsansätze systematisch bewertet, darunter eine Portal- bzw. \gls{federation} über Embedding, eine zentrale \gls{bi}-Aggregation mit Datenanbindung und neu aufgebauten Visualisierungen, Low-Code/No-Code-Plattformen für interne Tools sowie eine kundenspezifische Webanwendung. Kernbeitrag ist ein transparentes Scoring-Modell, das die Optionen anhand nachvollziehbarer Kriterien vergleichbar macht und in einer systematisch hergeleiteten Entscheidungsvorlage mündet; ergänzend wird die Robustheit der Empfehlung über eine Sensitivitätsanalyse der Kriteriengewichtung überprüft.

    Zur artefaktbasierten Evaluation wird ein \gls{mvp} umgesetzt, um zentrale Machbarkeitsannahmen und technische Risiken der favorisierten Alternative (insbesondere hinsichtlich Zugriffskontrolle, Datenabruf und Performance) anhand von mindestens zwei Quellsystemen und ausgewählten Kernmetriken zu verifizieren. Damit liefert die Arbeit eine begründete Entscheidungsvorlage einschließlich einer groben Umsetzungsroadmap sowie einen fokussierten \gls{poc}, der Risiken früh sichtbar macht und als Ausgangspunkt für eine spätere Skalierung auf weitere Datenquellen dienen kann.
}
\hypersetup{
    hidelinks
}
\begin{document}
    \begin{basimple}[
        img=images/logo.png,
        course=Wirtschaftsinformatik,
        title=Ein Mehrkriterien-Entscheidungsmodell zur Auswahl eines KPI-Hub-Integrationsansatzes in einer fragmentierten Enterprise-Systemlandschaft ,
        number=3005149,
        corrector={Janine Steidelmüller,Prof. Dr.-Ing. Jürgen Sachse},
        themedate=23. März 2026,
        returndate=\today,
        signature=images/signature.png,
        location=Dresden,
        type=thesis,
        % blockuntil=10. Juli 2029,  % Befristung Sperrvermerk (3 Jahre ab Abgabe) — einkommentieren wenn gewünscht
        kidoc=ki-genehmigung.pdf,  % KI-Genehmigung — vollständig unterschrieben (Sachse 21.04., Janine + Haakon 24.04.)
    ]

        \section{Einleitung}

\subsection{Ausgangssituation und Problemstellung}
% Hänel: Herleitung der Problemstellung — Warum ist das ein Problem?
% Bezug zum Unternehmensumfeld (SAP Signavio P&E), Branche, Allgemeinheit
% Nachweis durch Studien, Statistiken, Forschungsergebnisse (keine internen Dokumente/Blogs)

\subsection{Lösungsrelevanz und Motivation}
% Hänel: Diskrepanz aktuell ↔ gewünscht
% Hindernisse / Schwierigkeiten
% Negative Konsequenzen (Ineffizienz, Fehleranfälligkeit, mangelnde Transparenz)
% Aktueller Handlungsbedarf

\subsection{Einordnung in die Wirtschaftsinformatik}
% Hänel: Einordnung in die Wirtschaftsinformatik
% Bezug zu IS-Forschungsfeldern (z.B. Enterprise Systems, BI, Data Governance)

\subsection{Zielsetzung und Forschungsfrage}
% Forschungsfrage der Arbeit präzise formulieren
% Abgrenzung: Was wird untersucht, was nicht?

\subsection{Aufbau der Arbeit}
% Kurzer Überblick über die Kapitelstruktur

        \section{Grundlagen und Stand der Forschung}
        \section{Methodik}

\subsection{Forschungsansatz}
% Hänel: Welcher Forschungsansatz wird warum gewählt?
% Design Science Research nach Heinzl + Riedl (Referenz einfügen)
% Begründung warum DSR für diese Arbeit geeignet ist

\subsection{Literaturrecherche}
% Hänel: Rechercheprozess dokumentieren (Suche, Auswertung, Analyse, Interpretation)
% Verwendete Datenbanken, Suchbegriffe, Ein-/Ausschlusskriterien
% Anzahl gefundener / verwendeter Quellen

\subsection{Ergebniserwartung}
% Hänel: Ergebniserwartung anhand der Methoden aufstellen
% Was soll aus der Methodik herauskommen?
% Wie wird Zielerreichung gemessen?

\subsection{Vorgehen im Überblick}
% Hänel: Wie erfolgt die Untersuchung?
% Überblick über die einzelnen Phasen (Anforderungsanalyse → Lösungsraum → Bewertungsmodell → MVP)

        \section{Ausgangssituation und Systemlandschaft}

\subsection{Organisatorischer Kontext}
% Hänel: Vor welchem konkreten Hintergrund tritt das Problem auf?
% SAP Signavio Product & Engineering — Bereichsbeschreibung (ohne vertrauliche Details)
% Nachweis durch öffentlich zugängliche Strukturbeschreibungen

\subsection{Bestehende Systemlandschaft}
% Hänel: Überblick über eingesetzte Informationssysteme
% SAP Analytics Cloud, JIRA, ServiceNow, PDF/Excel-Reports
% Kommunikationsbeziehungen und Datenflüsse zwischen den Systemen

\subsection{Ist-Zustand der KPI-Versorgung}
% Wie werden KPIs heute erhoben und bereitgestellt?
% Manuelle Prozesse, Fragmentierung, Medienbrüche

\subsection{Datenbereitstellung und Schnittstellenlandschaft}
% Janine (21.04.): Thema Datenbereitstellung muss explizit abgedeckt werden
% Welche APIs, Exportformate, Datenflüsse existieren zwischen den Systemen?
% Technische Integrationsmöglichkeiten je Quellsystem (REST, ODATA, CSV, Embedding)
% Authentifizierungs- und Autorisierungsmechanismen (OAuth, SAML, API-Keys)
% Lizenzrechtliche und Security-Einschränkungen

\subsection{Identifizierte Problemfelder}
% Hänel: Nachweis durch Dokumentation von Strukturen, Problemen, Herausforderungen
% Konkrete Pain Points: Zeitaufwand, Fehleranfälligkeit, fehlende systemübergreifende Sicht

        \section{Anforderungsanalyse}

\subsection{Vorgehen bei der Anforderungserhebung}
% Methode: Stakeholder-Interviews, Use-Case-Analyse
% Auswahl der Stakeholder und Begründung

\subsection{Stakeholder und Use Cases}
% Zentrale Nutzergruppen und deren Informationsbedürfnisse
% Priorisierte Use Cases

\subsection{Funktionale Anforderungen}
% Datenquellen, Metriken, Aktualisierungsintervalle, Zielgruppen

\subsection{Nicht-funktionale Anforderungen}
% Performance, Datenschutz, RBAC, Compliance, Wartbarkeit, Lizenzkosten

\subsection{Priorisierung der Anforderungen}
% MoSCoW oder vergleichbare Methode

        \section{Lösungsraum und Architektur-/Tooloptionen}

\subsection{Identifikation relevanter Lösungsansätze}
% Systematische Herleitung der Kandidaten aus den Anforderungen

\subsection{Portal-/Dashboard-Föderation (Embedding)}
% Beschreibung, technische Voraussetzungen, Einschränkungen

\subsection{Zentrale BI-Aggregation}
% Beschreibung, Datenanbindung, neu aufgebaute Visualisierungen

\subsection{Low-Code / No-Code Plattformen}
% Interne Tools, Citizen-Developer-Ansatz

\subsection{Kundenspezifische Webanwendung}
% Custom Development, maximale Flexibilität, höherer Aufwand

\subsection{Vergleichende Übersicht}
% Tabellarischer Überblick aller Optionen vor der Bewertung

        \section{Entscheidungs- und Bewertungsmodell}

\subsection{Aufbau des Scoring-Modells}
% Kriterienauswahl und -begründung (aus Anforderungsanalyse abgeleitet)
% Gewichtung der Kriterien

\subsection{Bewertung der Lösungsoptionen}
% Systematische Bewertung jeder Option anhand des Scoring-Modells
% Transparente Herleitung der Punktevergabe

\subsection{Ergebnis und Empfehlung}
% Gesamtscoring, favorisierte Option
% Begründung der Entscheidung

\subsection{Sensitivitätsanalyse}
% Robustheitsprüfung: Wie verändert sich das Ergebnis bei veränderter Kriteriengewichtung?
% Hänel: Nachvollziehbare Kriterien → transparente Entscheidungsvorlage

        \section{MVP: Design, Umsetzung und Validierung}

\subsection{MVP-Zielsetzung und Scope}
% Hänel: Durchführung der Untersuchung
% Welche Machbarkeitsannahmen und technischen Risiken sollen verifiziert werden?
% Mindestens 2 Quellsysteme, ausgewählte Kernmetriken

\subsection{Architektur und Design}
% Technische Architektur des MVP
% Designentscheidungen und deren Begründung

\subsection{Umsetzung}
% Implementierungsschritte
% Eingesetzte Tools, Frameworks, APIs

\subsection{Validierung}
% Hänel: Ergebnisse präsentieren und Bedeutung in Bezug zur Problemstellung erläutern
% Test der Kernhypothesen (RBAC, Datenabruf, Performance)
% Dokumentation der Testergebnisse

        \section{Ergebnisse, Diskussion}

\subsection{Ergebnisse der Bewertung}
% Hänel: Ergebnisse präsentieren und deren Bedeutung in Bezug zur Problemstellung erläutern
% Zusammenfassung Scoring-Modell + MVP-Validierung

\subsection{Beantwortung der Forschungsfrage}
% Roter Faden: Bezug zur Forschungsfrage aus Kapitel 1

\subsection{Kritische Diskussion}
% Limitationen des Bewertungsmodells
% Grenzen der MVP-Validierung
% Generalisierbarkeit der Ergebnisse

\subsection{Vergleich mit Ergebniserwartung}
% Hänel: Ist das herausgekommen, was aus der Methodik erwartet wurde?

        \section{Fazit und Ausblick}

\subsection{Fazit}
% Hänel: Was bedeuten diese Ergebnisse in Bezug auf die Problemstellung?
% Warum das? Warum jetzt? Was ist neu? Für wen ist das interessant?
% Zielgruppe der Arbeit benennen

\subsection{Beitrag der Arbeit}
% Hänel: Und nun? Gut gemacht?
% Wissenschaftlicher und praktischer Beitrag zusammenfassen

\subsection{Ausblick}
% Hänel: Was muss im Weiteren getan werden, um die Forschung fortzusetzen?
% Skalierung auf weitere Quellsysteme
% Offene Fragen für Folgearbeiten

        \clearpage

        % appendix
        \begin{baappx}
            \section{Suchstrings der systematischen Literaturrecherche}
\label{apx:suchstrings}

Die nachfolgenden Tabellen dokumentieren die vollständigen Suchstrings für alle fünf thematischen Blöcke der systematischen Literaturrecherche (vgl.\ Abschnitt~\ref{sec:literaturrecherche}).
Die Strings sind für die primären Datenbanken \textit{ACM Digital Library} und \textit{IEEE Xplore} optimiert; für \textit{Springer Link} wurden längere Strings vereinfacht.
Notation: \texttt{AND}, \texttt{OR}, \texttt{*} (Trunkierung), \texttt{"..."} (Phrase).

% -----------------------------------------------------------------
% T1
% -----------------------------------------------------------------
\begin{batable}[label=tab:suchstrings-t1]{Suchstrings T1: Enterprise \gls{bi}-Integration und \gls{kpi}-Management}
\begin{tabular}{p{0.88\textwidth}}
\toprule
\textbf{Suchstring} \\
\midrule
\texttt{("KPI" OR "key performance indicator" OR "business intelligence") AND} \\
\texttt{("enterprise integration" OR "fragmented" OR "dashboard consolidation" OR "BI landscape")} \\[4pt]
\texttt{("KPI management" OR "KPI hub" OR "KPI dashboard") AND} \\
\texttt{("enterprise system" OR "information system" OR "SAP")} \\[4pt]
\texttt{"business intelligence fragmentation" OR "BI fragmentation"} \\[4pt]
\texttt{("KPI" OR "performance measurement") AND "enterprise information system*"} \\[4pt]
\textit{Deutsch (Springer~/ WISO):} \\
\texttt{("KPI" OR "Kennzahl") AND ("Enterprise" OR "Unternehmens-BI") AND ("Integration" OR "Fragmentierung")} \\
\bottomrule
\end{tabular}
\end{batable}

\noindent\textbf{Datenbanken T1:} ACM Digital Library, IEEE Xplore, Springer Link, EBSCO Business Source Complete

% -----------------------------------------------------------------
% T2
% -----------------------------------------------------------------
\begin{batable}[label=tab:suchstrings-t2]{Suchstrings T2: Integrationsparadigmen (\gls{et}, \gls{eip}, \gls{rest}, Embedded Analytics)}
\begin{tabular}{p{0.88\textwidth}}
\toprule
\textbf{Suchstring} \\
\midrule
\texttt{("ETL" OR "extract transform load") AND ("business intelligence" OR "data warehouse" OR "enterprise")} \\[4pt]
\texttt{("enterprise application integration" OR "EAI") AND ("integration pattern*" OR "message broker")} \\[4pt]
\texttt{("REST API" OR "RESTful") AND ("enterprise integration" OR "data integration")} \\[4pt]
\texttt{"embedded analytics" AND ("enterprise system*" OR "SAP" OR "ERP")} \\[4pt]
\texttt{("data mesh" OR "data fabric") AND "enterprise analytics"} \\[4pt]
\texttt{("integration paradigm*" OR "integration approach*") AND} \\
\texttt{("enterprise system*" OR "ERP" OR "business intelligence")} \\
\bottomrule
\end{tabular}
\end{batable}

\noindent\textbf{Datenbanken T2:} ACM Digital Library, IEEE Xplore, Springer Link

% -----------------------------------------------------------------
% T3
% -----------------------------------------------------------------
\begin{batable}[label=tab:suchstrings-t3]{Suchstrings T3: \gls{rbac} in Enterprise-Systemen}
\begin{tabular}{p{0.88\textwidth}}
\toprule
\textbf{Suchstring} \\
\midrule
\texttt{("role-based access control" OR "RBAC") AND} \\
\texttt{("business intelligence" OR "analytics" OR "dashboard" OR "enterprise")} \\[4pt]
\texttt{("RBAC" OR "access control") AND ("data governance" OR "information governance")} \\[4pt]
\texttt{("role-based access control" OR "RBAC") AND ("ERP" OR "SAP" OR "enterprise system*")} \\[4pt]
\texttt{"attribute-based access control" AND "enterprise analytics"} \\
\bottomrule
\end{tabular}
\end{batable}

\noindent\textbf{Datenbanken T3:} ACM Digital Library, IEEE Xplore, Springer Link

% -----------------------------------------------------------------
% T4
% -----------------------------------------------------------------
\begin{batable}[label=tab:suchstrings-t4]{Suchstrings T4: Mehrkriterien-Entscheidungsmethoden (\gls{ahp}, \gls{topsis}, Nutzwertanalyse)}
\begin{tabular}{p{0.88\textwidth}}
\toprule
\textbf{Suchstring} \\
\midrule
\texttt{("analytic hierarchy process" OR "AHP") AND} \\
\texttt{("software selection" OR "tool selection" OR "IT selection" OR "enterprise system")} \\[4pt]
\texttt{("TOPSIS" OR "technique for order preference") AND} \\
\texttt{("information system*" OR "software engineering" OR "IT evaluation")} \\[4pt]
\texttt{("multi-criteria decision" OR "MCDM" OR "multiple criteria") AND} \\
\texttt{("enterprise system selection" OR "IT architecture" OR "integration approach")} \\[4pt]
\texttt{"Nutzwertanalyse" AND ("Softwareauswahl" OR "IT-Entscheidung")} \\[4pt]
\texttt{("weighted scoring" OR "scoring model") AND} \\
\texttt{("information system*" OR "enterprise software" OR "tool selection")} \\
\bottomrule
\end{tabular}
\end{batable}

\noindent\textbf{Datenbanken T4:} ACM Digital Library, IEEE Xplore, Springer Link, EBSCO Business Source Complete

% -----------------------------------------------------------------
% T5
% -----------------------------------------------------------------
\begin{batable}[label=tab:suchstrings-t5]{Suchstrings T5: \gls{dsr}-Methodik}
\begin{tabular}{p{0.88\textwidth}}
\toprule
\textbf{Suchstring} \\
\midrule
\texttt{"design science research" AND "information system*"} \\[4pt]
\texttt{("DSR" OR "design science") AND ("evaluation" OR "artifact" OR "artefact")} \\
\bottomrule
\end{tabular}
\end{batable}

\noindent\textbf{Datenbanken T5:} ACM Digital Library, IEEE Xplore, Springer Link

        \end{baappx}
    \end{basimple}
\end{document}
